\section{Doubling and the first activation}\label{sec:activation}

Theorem~\ref{thm:bounded-transients} bounds the total number of reflection
classes which do not come from the whole loop space. We now compare that
number at a cut and its double. Equality in this comparison will require
each visible odd-stage bottom class to have a visible evaluation-cup
partner. At the first appearance of such a bottom, local evaluation
vanishing prevents its partner from appearing. This gives the final
contradiction.

Retain the finite-prime assumption and write
\[
 V(a)=\sum_j\rank(B_\infty^j\longrightarrow B_a^j),
 \qquad
 T(a)=\sum_j\dim B_a^j-V(a).
\]
Both sums are finite at a regular cut. We first record the full-group
rank calculation that makes the doubling comparison exact.

\begin{lemma}[Free rank and torsion at a finite cut]\label{act:module-rank}
Suppose that $2L>0$ is regular, and put
$H=\bigoplus_j H_G^j(X_{2L},M)$. This is a finitely generated graded
$\kk[U]$-module. Its torsion submodule $\mathcal T$ is finite-dimensional
and annihilated by a power of $U$. If
\[
 t=\dim(\mathcal T/U\mathcal T)=\dim\ker(U|_H),
 \qquad \rho=\rank_{\kk[U]}(H/\mathcal T),
\]
then
\begin{equation}\label{act:total-rank}
 \sum_j\dim B_{2L}^j=\rho+2t,
 \qquad \rho=\sum_j\dim B_L^j.
\end{equation}
\end{lemma}

\begin{proof}
The standard-plane Gysin sequence in Proposition~\ref{prop:gysin}
embeds $H/UH$ into $\bigoplus_j B_{2L}^j$, which has finite dimension.
Lift a finite homogeneous basis of $H/UH$. These lifts generate $H$:
subtract their contributions modulo $U$, divide the remainder by $U$,
and induct down the nonnegative degree. This is an induction decreasing
degree by two at each step.

If a polynomial $p(U)$ has nonzero constant term, it acts injectively
on $H$. Indeed, the lowest nonzero homogeneous component of a nonzero
element survives in its product with $p(U)$. Factoring an annihilating
polynomial as $U^s p(U)$ therefore shows that every torsion element is
killed by a power of $U$. The ring $\kk[U]$ is Noetherian, so $\mathcal T$
has finitely many generators and is killed by one common power of $U$.
It is consequently finite-dimensional. Rank-nullity on $\mathcal T$
gives
\[
 \dim\ker(U|_H)=\dim\ker(U|_{\mathcal T})
 =\dim(\mathcal T/U\mathcal T).
\]
Moreover $UH\cap\mathcal T=U\mathcal T$, since $Uh$ torsion implies
that $h$ is torsion.

Set $\mathcal F=H/\mathcal T$. Multiplication by $U$ is injective on
$\mathcal F$. Lift a homogeneous basis of $\mathcal F/U\mathcal F$.
Degree induction again proves generation. These lifts are free: in a
homogeneous relation cancel the least power of $U$ and reduce modulo
$U$. Independence of the selected basis excludes the terms with that
least power, a contradiction. Thus $\mathcal F$ has a homogeneous free
basis of size $\rho$, and
\[
 \dim(H/UH)=\rho+t.
\]
Summing the standard-plane Gysin sequence over degrees gives the first
identity in \eqref{act:total-rank}, since its two contributions are
$H/UH$ and a shifted copy of $\ker U$.

For sufficiently large $D$, the torsion no longer contributes in
degrees $D,D-1$. The dimensions of those two free degrees add to
$\rho$, one contribution from each free generator according to parity.
The weight-two-plane Gysin sequence gives
\[
 \rho=\dim H_G^D(X_{2L},M)+\dim H_G^{D-1}(X_{2L},M)
 =\dim H_K^D(X_{2L},M).
\]
By Lemma~\ref{bud:localization} and \eqref{bud:fixed-product}, the last
dimension equals $\sum_j\dim B_L^j$ in high degree. This proves the
second identity without requiring surjectivity of full-group restriction
to the fixed locus.
\end{proof}

Use the global basis from Theorem~\ref{thm:global-generators}:
\[
 \xi_{m,k}=w^kX_m,\qquad \eta_{m,k}=e\xi_{m,k},
 \qquad m\ge1,\quad 0\le k\le d.
\]
At a regular cut $a$, let $b_{m,k}(a)$ be one if $\xi_{m,k}$ restricts
nontrivially and zero otherwise. Let $h_{m,k}(a)$ be one if the
restriction of $\eta_{m,k}$ is independent modulo all bottom
restrictions in its degree, and zero otherwise. Degrees within the
bottom family are distinct, as are degrees within the top family:
their consecutive starting degrees differ by $2d$, while each family
at a stage spans an interval of width $d$. Therefore
\begin{equation}\label{act:visibility-count}
 V(a)=\sum_{m,k}\bigl(b_{m,k}(a)+h_{m,k}(a)\bigr),
 \qquad h_{m,k}(a)\le b_{m,k}(a).
\end{equation}
The inequality follows from $\eta_{m,k}=e\xi_{m,k}$. Write $V_o(a)$
and $V_e(a)$ for the contributions of odd and even stages $m$ to this
sum. Stage parity here refers to the round-filtration label $m$.

\begin{theorem}[Doubling and its equality condition]\label{thm:doubling}
For every regular $2L>0$,
\begin{equation}\label{act:doubling}
 T(2L)\ge T(L).
\end{equation}
If equality holds, then
\begin{equation}\label{act:odd-equality}
 h_{m,k}(2L)=b_{m,k}(2L)
 \qquad(m\ \text{odd},\ 0\le k\le d).
\end{equation}
\end{theorem}

\begin{proof}
Use the notation of Lemma~\ref{act:module-rank}, and put
\[
 N=\sum_{m\ \mathrm{odd},k} b_{m,k}(2L).
\]
The selected odd bottoms have global $G$-equivariant $U$-torsion lifts
by Theorem~\ref{thm:global-generators}. Their restrictions to $H/UH$
are independent, since that quotient injects into reflection cohomology
and their bottom restrictions are independent. They thus occupy an
$N$-dimensional subspace of $\mathcal T/U\mathcal T$. Write
\begin{equation}\label{act:residual-torsion}
 t=N+J_0,\qquad J_0\ge0.
\end{equation}

The selected even bottoms and tops also have global full-group lifts.
Their classes modulo $U$ are jointly independent with the selected odd
bottoms: this follows from their reflection restrictions and the rule
that a top is counted only modulo all bottoms in its degree. In the
projection
\[
 H/UH\longrightarrow\mathcal F/U\mathcal F,
\]
the kernel is $\mathcal T/U\mathcal T$, of dimension $N+J_0$.
The odd span already occupies $N$ of these dimensions. Hence the image
of the even span has dimension at least $V_e(2L)-J_0$.

Choose a maximal homogeneous independent family among these even
projections; its size is at least $V_e(2L)-J_0$.
Multiply the corresponding global lifts by powers of $U$ to put them in
two common large degrees $D,D-1$, according to parity. Independence persists in
$\mathcal F$: from a relation in one of the common degrees, cancel
the least power of $U$ and reduce modulo $U$, then repeat if necessary.
Thus the sum of the ambient full-group global-image ranks in these
two degrees is at least $V_e(2L)-J_0$.

We compare these images with the fixed locus, including the image-level
Gysin maps. The subspace $X_{2L}^\tau$ is $G$-invariant. Under its
identification with $X_L$, rotations act with doubled parameter and
reflection acts as before. In sufficiently high degree the map
\[
 H_G^q(X_{2L},M)\longrightarrow H_G^q(X_{2L}^\tau,M)
\]
is injective. Indeed, compose it with the injective fixed $G$-to-$K$
restriction. The same composite factors through the ambient
$G$-to-$K$ injection and the high actual $K$ restriction isomorphism
of Lemma~\ref{bud:localization}.

For $H'=G,K$, set
\[
 I_{H'}^q=\im\bigl[H_{H'}^q(\Lambda^\tau,M)
       \longrightarrow H_{H'}^q(X_{2L}^\tau,M)\bigr].
\]
Naturality places the ambient global images inside these fixed global
images. Apply the weight-two-plane Gysin sequence to the fixed global
and fixed finite pairs. Its global push is surjective, so on images
the finite push gives a surjection
\[
 \pi:I_K^D\longrightarrow I_G^{D-1}.
\]
The finite $G$-to-$K$ injection identifies $I_G^D$ with a subspace of
$\ker\pi$. Consequently
\[
 \dim I_G^D+\dim I_G^{D-1}\le\dim I_K^D.
\]
On the fixed $K$-spaces the global restriction is the coefficientwise
map $B_\infty[v]\to B_L[v]$. Since $B_L$ has finite total dimension,
$\dim I_K^D=V(L)$ for sufficiently large $D$. Choose $D$ large enough
for all the preceding comparisons. We have proved
\begin{equation}\label{act:image-inequality}
 V(L)\ge V_e(2L)-J_0.
\end{equation}

Subtract global-image dimensions from \eqref{act:total-rank} and use
\eqref{act:residual-torsion}--\eqref{act:image-inequality}:
\begin{align*}
 T(2L)-T(L)
 &=2(N+J_0)-V_o(2L)-V_e(2L)+V(L)\\
 &\ge J_0+2N-V_o(2L)\\
 &=J_0+\sum_{m\ \mathrm{odd},k}
       \bigl(b_{m,k}(2L)-h_{m,k}(2L)\bigr)\ge0.
\end{align*}
All summands in the final line are nonnegative. Equality therefore
forces \eqref{act:odd-equality} for each individual odd bottom.
\end{proof}

\begin{proof}[Completion of the proof of Theorem~\ref{thm:spheres} for $n\ge3$]
Suppose that there are only finitely many prime geodesic images.
By Theorem~\ref{thm:bounded-transients}, $T$ is a bounded nonnegative
integer on regular positive cuts. It is constant on every no-critical
interval, so its maximum is attained on a nonempty bounded open regular
interval $I\subset(0,\infty)$.

Every regular point $x\in2^kI$ has regular successive halves: a critical
half would double to a critical point. Iterating \eqref{act:doubling}
therefore gives $T(x)\ge T(x/2^k)$, which is already the global maximum.
Thus every regular point in every expanded interval attains that maximum.
For $k\ge1$, equality holds between $x$ and $x/2$, so
\eqref{act:odd-equality} holds at every such $x$. This uses monotonicity
under doubling only; no monotonicity between arbitrary cuts is asserted.

Take $k$ large and choose regular $x_1<x_2$ in the lower and upper
quarters of $2^kI$. The interval of degrees strictly between
$\alpha x_1$ and $\alpha x_2$, after removing the fixed endpoint margins
of Lemma~\ref{trans:degree-gaps}, has length tending to infinity.
Remove also the fixed-width bands about every $\beta_i x_1$ and
$\beta_i x_2$. They remove a bounded total number of possible degrees.
The odd bottom degrees
\[
 a_m=(2m-1)d,\qquad m\ \text{odd},
\]
form an arithmetic progression of spacing $4d$. For sufficiently large
$k$, one such degree $j=a_m$ remains, beyond the fixed low-degree range.
Lemma~\ref{trans:degree-gaps} makes $\xi_{m,0}$ invisible at $x_1$
and visible at $x_2$. At $x_1$ the degree lies in an absent gap, where
the global image is zero, or above the last band, where the sublevel
group is zero. At $x_2$ it lies in a present gap where global reflection
restriction is injective.

A global class nonzero at a smaller sublevel cannot become zero at a
larger one, since restriction to the smaller sublevel factors through
the larger. There are only finitely many events between $x_1$ and $x_2$.
Choose the first event at which $\xi_{m,0}$ becomes nonzero, and tight
regular cuts $x_-<x_+$ around it, still inside $2^kI$. The exact sequence
of the crossing puts its post-event restriction in the image of
\[
 H_R^j(X_{x_+},X_{x_-})\longrightarrow B_{x_+}^j,
\]
because its pre-event restriction is zero. By
Theorem~\ref{thm:local-control}, cup with $e$ vanishes on this whole
reflection crossing, including simultaneous prime labels. Naturality
of the cup action therefore gives
\[
 \eta_{m,0}|_{X_{x_+}}=e\,\xi_{m,0}|_{X_{x_+}}=0.
\]
Hence $b_{m,0}(x_+)=1$ while $h_{m,0}(x_+)=0$. Since $m$ is odd,
this contradicts \eqref{act:odd-equality} at the regular point $x_+$
of the expanded interval. The finite-prime-image assumption is impossible.
\end{proof}
